\begin{scanpage}{0064}{47}
The vector potential of the currents flowing in the indicated sector is determined by the equation
\[\left.\begin{aligned}A_y&=-\frac1{2\pi}H_{0x}\frac{e^{ikR}}R\int_{-\Phi_0}^{\Phi_0}d\psi\int_0^a e^{-ik\rho\sin\theta\cos\psi}\rho\,d\rho,\\A_x&=A_z=0.\end{aligned}\right\}\tag{7.17}\]
Taking into account the condition $ka\sin\theta\gg1$, one may show that the field created by the currents of this sector will equal
\[\left.\begin{aligned}E_\theta=H_\phi&\sim H_{0x}\sqrt{\frac a{\pi k\sin\theta}}\cos\theta\times\\&\quad\times\frac{\sin\phi}{\sin\theta}\frac{e^{ikR}}R e^{-ika\sin\theta+i\frac\pi{5}}+O\left(\frac1{\sqrt{ka\sin\theta}}\right),\\E_\phi=-H_\theta&\sim H_{0x}\sqrt{\frac a{\pi k\sin\theta}}\times\\&\quad\times\frac{\cos\phi}{\sin\theta}\frac{e^{ikR}}R e^{-ika\sin\theta+i\frac\pi{5}}+O\left(\frac1{\sqrt{ka\sin\theta}}\right).\end{aligned}\right\}\tag{7.18}\]
The amplitude of the expressions which have been found is approximately $\sqrt2$ times larger than the amplitude of the first terms in Equation (7.13). Moreover, expressions (7.18) and the corresponding terms in Equation (7.13) differ slightly in their phases: the first have the factor $e^{i\frac\pi{5}}$, and the latter --- the factor $-e^{i\frac\pi4}$. The result obtained is similar to the well-known thesis in optics that the effect of a wave is equal to the effect of half of the first Fresnel zone (see, for example [27], p.\ 132).

In the vicinity of the directions $\theta=0$ and $\theta=\pi$, when the azimuthal components lose their meaning, for the purpose of studying the fringing field it is more convenient to use the Cartesian components
\[\left.\begin{aligned}E_x&=(E_\theta\cos\theta+E_R\sin\theta)\cos\phi-E_\phi\sin\phi,\\E_y&=(E_\theta\cos\theta+E_R\sin\theta)\sin\phi+E_\phi\cos\phi.\end{aligned}\right\}\tag{7.19}\]
Turning to Equations (7.11), we find that when $\theta=0$ and $\theta=\pi$
\[E_x=0,\quad E_y=-iH_{0x}\frac{ka^2}2\frac{e^{ikR}}R.\tag{7.20}\]
\end{scanpage}
